\newpage
\section{Perfil de referência NACA 1411}

O perfil de partida adotado é o NACA 1411, conforme o roteiro. Como a asa
exige, na estação de referência, a espessura relativa mínima de
$(t/c)_{\mathrm{ref}} = 0{,}1772$ no plano normal
(Tabela~\ref{tab:ref_section}), o perfil foi reescalado antes de qualquer
análise, multiplicando a distribuição de espessura pelo fator
$0{,}1772/0{,}11$ e preservando a linha de arqueamento. Reescalar as
ordenadas completas inflaria também o arqueamento e mudaria a sustentação
do perfil de partida. Sem esse reescalonamento, partida e otimizado não
seriam comparáveis, por terem espessuras estruturais diferentes.

O procedimento para encontrar o ângulo de ataque de projeto foi o seguinte.
Fixado $M_n = 0{,}7196$ no \texttt{single\_run}, o solver Euler foi executado
na mesma malha usada na otimização (malha-O de $31\times 49$ pontos) e a
variável \texttt{alpha} foi ajustada pelo método de Newton, usando a derivada
$\partial c_\ell/\partial \alpha$ fornecida pelo método adjunto, até que o
$c_\ell$ retornado igualasse o de projeto com tolerância de $10^{-3}$. A
convergência exigiu apenas 3 avaliações, partindo de $\alpha = 3^\circ$
(onde o perfil sustenta $c_\ell = 0{,}667$).

O ângulo encontrado foi $\alpha = 3{,}567^\circ$, com $c_\ell = 0{,}7319$,
confirmando o casamento com o $c_{\ell,\mathrm{ref}}$ da
Tabela~\ref{tab:ref_section}. Nessa condição o perfil de partida apresenta
$c_d = 0{,}06705$ e $c_m = -0{,}0680$. O arrasto é elevado, dominado pelo
arrasto de onda de um choque forte no extradorso, como será visto na
Seção~5. É esse o número que a otimização da Seção~4 ataca.
